\section{Problem Statement}

\subsection{Problem Definition}

This work checks whether a \textit{fully decentralized} LoRa-based V2V mesh can provide reliable and efficient
multi-hop safety-message dissemination when there is no supporting infrastructure.

This is difficult because LoRa links impose strict communication constraints, including low data rate, long
Time-on-Air, duty-cycle limits, and shared-medium contention \cite{lora_mesh_iot,routing_mesh_2020}. If the
protocol is not controlled carefully, these constraints can cause collisions, redundant rebroadcasts, and
unstable dissemination in dynamic vehicular topologies \cite{vdtn_routing_eval,lora_mesh_emergency}. So the
central problem is designing a decentralized protocol stack that still reaches far enough, but keeps channel
usage bounded and the behavior predictable.

\subsection{Objectives and Goals}

The project objectives are:

\begin{itemize}
	\item Implement a three-layer architecture with clear separation between application logic, forwarding control,
	and radio-facing transmission services.
	\item Implement decentralized network-layer dissemination using managed flooding, duplicate suppression,
	hop-bounded propagation, and policy-aware relay timing.
	\item Implement link-layer channel-aware behavior, including sensing-driven transmission admission and
	optional reliability signaling for critical traffic classes.
	\item Demonstrate that the system can deliver safety messages across multiple hops with bounded redundancy,
	without dependence on RSUs, fixed gateways, or centralized coordination.
\end{itemize}

\subsection{Significance of the Problem}

This problem matters for a few reasons. For a final-year project, it gives a low-cost and reproducible way to
study decentralized vehicular networking. It also provides practical evidence of infrastructure-free
dissemination behavior for scenarios where fixed communication support may be sparse or temporarily
unavailable \cite{lora_mesh_emergency,lora_reliable_v2v_2025}. Finally, it shows a workflow where the same
network-layer logic can be tested in simulation and then used in embedded deployment, which helps build
confidence in the protocol behavior before trying larger real-world setups.

